\chapter{FORMS, ORIENTATION, AND INTEGRATION}

\section{Differential Forms}
\textit{October 5th.}

We translate our theory of differential forms on $\R^{n}$ onto smooth manifolds. A 1-form on a manifold $M$ is section of the cotangent bundle $T^{\ast}M$. That is, a 1-form $\omega$ on $M$ is a map $\omega:M \to T^{\ast}M = \bigsqcup_{p \in M} T^{\ast}_{p}M$ such that $\omega(p) \in T^{\ast}_{p}M$ for all $p \in M$. This is analagous to a vector field $X$ which is a map $X:M \to TM = \bigsqcup_{p \in M} T_{p}M$ such that $X(p) \in T_{p}M$ for all $p \in M$. Given a smooth function $f:M \to \R$, we have the differential $df:M \to T^{\ast}M$ defined as $df(p) = df_{p}$, where $df_{p}:T_{p}M \to \R$ is the linear map defined as $df_{p}(X_{p}) = X_{p}(f)$ for all $X_{p} \in T_{p}M$. In particular, we have that $df$ is a 1-form on $M$.

Moreover, given a chart $(U,\varphi) = (U,x^{1},\ldots,x^{n})$ on $M$, the $x^{i}$'s can, again, be thought of as coordinate functions. With $\{\frac{\partial}{\partial x^{i}}\}_{i=1}^{n}$ the standard basis for $T_{p}M$, $p$ being in $U$, the partials $\{(dx^{i})_{p}\}$ forms the respective dual basis for $T_{p}^{\ast}M$. Thus, any $\omega_{p} \in T_{p}^{\ast}M$ can be written as
\begin{align}
    \omega_{p} = \sum_{i=1}^{n} a_{i}(p)(dx^{i})_{p}
\end{align}
where $a_{i}(p)$'s are real numbers. More generally, any 1-form $\omega$ on $U$ can be written as
\begin{align}
    \omega = \sum_{i=1}^{n} a_{i} dx^{i}
\end{align}
where $a_{i}:U \to \R$ are real-valued functions on $U$.

\begin{definition}
    The \eax{cotangent bundle} of a smooth manifold $M$ is the vector bundle $T^{\ast}M = \bigsqcup_{p \in M} T^{\ast}_{p}M$ over $M$, with projection $\pi:T^{\ast}M \to M$ defined as $\pi(\omega_{p}) = p$ for all $\omega_{p} \in T^{\ast}_{p}M$. A \eax{1-form} on $M$ is a section of the cotangent bundle, that is, a map $\omega:M \to T^{\ast}M$ such that $\omega(p) \in T^{\ast}_{p}M$ for all $p \in M$.
\end{definition}

We can, again, put a topology on $T^{\ast}M$ by using the charts on $M$, just as we did with the tangent bundle $TM$. Let $(U,\varphi) = (U,x^{1},\ldots,x^{n})$ be a chart on $M$. Also adopt the notation $T^{\ast}U = \pi^{-1}(U) = \bigsqcup_{p \in U} T^{\ast}_{p}M$. Then we have the bijective map
\begin{align}
    \tilde{\varphi}:T^{\ast}U \to \varphi(U) \times \R^{n}, \qquad \omega_{p} \mapsto (\varphi(p),a_{1}(p),\ldots,a_{n}(p)).
\end{align}
With the standard topology on $\varphi(U) \times \R^{n}$, we can give $T^{\ast}U$ the topology induced by $\tilde{\varphi}$. Then, letting $\cB_{U}$ be the collection of all open sets in $T^{\ast}U$, we can give $T^{\ast}M$ the topology generated by $\cB = \bigcup_{(U,\varphi)} \cB_{U}$, where the union is taken over all charts $(U,\varphi)$ on $M$. The exact same analysis as tangent bundles shows that $T^{\ast}M$ is a smooth manifold of dimension $2n$, and the projection $\pi:T^{\ast}M \to M$ is a smooth map.

A 1-form $\omega$ on $M$ is said to be smooth if $\omega$ is a smooth section of $\pi$. That is, $\omega$ is smooth if $\omega:M \to T^{\ast}M$ is a smooth map between smooth manifolds.

\begin{proposition}
    Let $\omega$ be a 1-form on a smooth manifold $M$. Then the following are equivalent.
    \begin{enumerate}
        \item $\omega$ is smooth.
        \item For any chart $(U,\varphi) = (U,x^{1},\ldots,x^{n})$ on $M$, the functions $a_{i}:U \to \R$ defined as $\omega = \sum_{i=1}^{n} a_{i} dx^{i}$ are smooth.
    \end{enumerate}
\end{proposition}

Let $X$ be a vector field on a smooth manifold $M$, and let $\omega$ be a 1-form on $M$. Then we can define the map $\omega(X):M \to \R$ as $\omega(X)(p) = \omega_{p}(X_{p})$ for all $p \in M$. That is, $\omega(X)$ is the function that takes a point $p \in M$ and returns the value of the linear map $\omega_{p}:T_{p}M \to \R$ evaluated at the vector $X_{p} \in T_{p}M$.

\begin{proposition}
    Let $\omega$ be a 1-form on a smooth manifold $M$. Then $\omega$ is a smooth 1-form if and only if $\omega(X)$ is a smooth function on $M$ for all smooth vector fields $X$ on $M$.
\end{proposition}
\begin{proof}
    Let $\omega$ be a smooth 1-form on $M$. Let $X$ be a smooth vector field on $M$. Then the map $X:M \to TM$ is smooth, and so the composition $\omega(X) = \omega \circ X:M \to \R$ is smooth. Conversely, suppose that $\omega(X)$ is a smooth function on $M$ for all smooth vector fields $X$ on $M$. Let $(U,\varphi) = (U,x^{1},\ldots,x^{n})$ be a chart on $M$. Then the partials $\{\frac{\partial}{\partial x^{i}}\}_{i=1}^{n}$ form a basis for the tangent space $T_{p}M$ for all $p \in U$, and so we can write
    \begin{align}
        \omega = \sum_{i=1}^{n} a_{i} dx^{i}
    \end{align}
    for some functions $a_{i}:U \to \R$. Then we have
    \begin{align}
        a_{i} = dx^{i}\left(\frac{\partial}{\partial x^{i}}\right) = \omega\left(\frac{\partial}{\partial x^{i}}\right)
    \end{align}
    for all $i = 1,\ldots,n$. Since $\frac{\partial}{\partial x^{i}}$ is a smooth vector field on $U$, we have that $\omega\left(\frac{\partial}{\partial x^{i}}\right)$ is a smooth function on $U$, and so each of the functions $a_{i}$ are smooth. Thus, $\omega$ is a smooth 1-form on $M$.
\end{proof}

Recall that $\mf{X}(M)$ is the real vector space of all smooth vector fields on a smooth manifold $M$, which is also a $C^{\infty}(M)$-module via $(fX)(p) = f(p)X_{p}$ for all $f \in C^{\infty}(M)$, $X \in \mf{X}(M)$, and $p \in M$. Similarly, we have the real vector space $\Omega^{1}(M)$ of all smooth 1-forms on $M$, which is also a $C^{\infty}(M)$-module via $(f\omega)(p) = f(p)\omega_{p}$ for all $f \in C^{\infty}(M)$, $\omega \in \Omega^{1}(M)$, and $p \in M$. Therefore, we can define the map $\Omega^{1}(M) \times \mf{X}(M) \to C^{\infty}(M)$ as $(\omega,X) \mapsto \omega(X)$, which is $\R$-bilinear and $C^{\infty}(M)$-bilinear. That is, we have
\begin{align}
    (\omega_{1}+\omega_{2})(X) = \omega_{1}(X) + \omega_{2}(X), \qquad \omega(X_{1}+X_{2}) = \omega(X_{1}) + \omega(X_{2}), \notag \\ (f\omega)(X) = f(\omega(X)), \qquad \omega(fX) = f(\omega(X))
\end{align}
Now, if $T:V \to W$ is a linear map between finite-dimensional vector spaces $V$ and $W$, then the dual map $T^{\ast}:W^{\ast} \to V^{\ast}$ is defined as $T^{\ast}(\omega) = \omega \circ T$ for all $\omega \in W^{\ast}$. That is, $T^{\ast}$ takes a linear functional $\omega:W \to \R$ and returns the linear functional $\omega \circ T:V \to \R$. The dual map is also linear, and satisfies $(S \circ T)^{\ast} = T^{\ast} \circ S^{\ast}$ for all linear maps $T:V \to W$ and $S:W \to U$ between finite-dimensional vector spaces $V$, $W$, and $U$. We can extend this notion to smooth maps between smooth manifolds.

Let $F:N \to M$ be a smooth map between smooth manifolds $N$ and $M$. Then the differential $F_{\ast,p}:T_{p}N \to T_{F(p)}M$ is a linear map between the tangent spaces at $p \in N$ and $F(p) \in M$. Therefore, we can define the dual map $F^{\ast}_{p}:T^{\ast}_{F(p)}M \to T^{\ast}_{p}N$ as $F^{\ast}_{p}(\omega) = \omega \circ F_{\ast,p}$ for all $\omega \in T^{\ast}_{F(p)}M$. That is, $F^{\ast}_{p}$ takes a linear functional $\omega_{F(p)}:T_{F(p)}M \to \R$ and returns the linear functional $\omega_{F(p)} \circ F_{\ast,p}:T_{p}N \to \R$. The dual map is also linear, and satisfies $(G \circ F)^{\ast}_{p} = F^{\ast}_{p} \circ G^{\ast}_{F(p)}$ for all smooth maps $F:N \to M$ and $G:M \to P$ between smooth manifolds $N$, $M$, and $P$. If $\omega$ is a 1-form on $M$, then we can define the pullback of $\omega$ by $F$ as the 1-form $F^{\ast}\omega$ on $N$ defined as $(F^{\ast}\omega)_{p} = F^{\ast}_{p}(\omega_{F(p)})$ for all $p \in N$.

\begin{proposition}
    Let $\omega$ be a smooth 1-form on a smooth manifold $M$, and let $F:N \to M$ be a smooth map between smooth manifolds $N$ and $M$. Then the pullback $F^{\ast}\omega$ is a smooth 1-form on $N$.
\end{proposition}
\begin{proof}
    Let $(U,\varphi) = (U,x^{1},\ldots,x^{n})$ be a chart on $M$, and let $(V,\psi) = (V,y^{1},\ldots,y^{m})$ be a chart on $N$. Then we have
    \begin{align}
        \omega = \sum_{i=1}^{n} a_{i} dx^{i}, \qquad F^{\ast}\omega = \sum_{j=1}^{m} b_{j} dy^{j}
    \end{align}
    for some smooth functions $a_{i}:U \to \R$ and $b_{j}:V \to \R$. Then we have
    \begin{align}
        b_{j}(p) = (F^{\ast}\omega)_{p}\left(\frac{\partial}{\partial y^{j}}\right) = F^{\ast}_{p}(\omega_{F(p)})\left(\frac{\partial}{\partial y^{j}}\right) = \omega_{F(p)}(F_{\ast,p}\left(\frac{\partial}{\partial y^{j}}\right)) = \sum_{i=1}^{n} a_{i}(F(p)) dx^{i}(F_{\ast,p}\left(\frac{\partial}{\partial y^{j}}\right)).
    \end{align}
    Since $F_{\ast,p}\left(\frac{\partial}{\partial y^{j}}\right)$ is a smooth vector field on $N$, we have that $dx^{i}(F_{\ast,p}\left(\frac{\partial}{\partial y^{j}}\right))$ is a smooth function on $N$. Therefore, each of the functions $b_{j}$ are smooth, and so $F^{\ast}\omega$ is a smooth 1-form on $N$.
\end{proof}

\subsection{$k$-forms}
\noindent \textit{October 7th.}

Moving on to $k$-forms, recall that $A_{k}(V)$ is the space of all alternating $k$-linear maps $\omega:V^{k} \to \R$ on a finite-dimensional vector space $V$. It is of dimension $\binom{n}{k}$, and if $\{\alpha_{1},\ldots,\alpha_{n}\}$ is a basis for $V^{\ast}$, then the wedge products $\{\alpha_{i_{1}} \wedge \cdots \wedge \alpha_{i_{k}} \mid 1 \leq i_{1} < \cdots < i_{k} \leq n\}$ form a basis for $A_{k}(V)$. We can extend this notion to smooth manifolds.

Let $M$ be a smooth manifold of dimension $n$, and let $p \in M$. We have
\begin{align}
    A_{k}(T_{p}M) \defeq \{\omega_{p}:T_{p}M^{k} \to \R \mid \omega_{p} \text{ is alternating and $k$-linear}\} = \bigwedge^{k}T_{p}^{\ast}M.
\end{align}
A $k$-form on $M$ is a map
\begin{align}
    \omega:M \to \bigsqcup_{p \in M} A_{k}(T_{p}M) = \bigsqcup_{p \in M} \bigwedge^{k}T_{p}^{\ast}M
\end{align}
such that $\omega(p) \in A_{k}(T_{p}M)$ for all $p \in M$. If we denote $\cI_{k,n} = \{(i_{1},\ldots,i_{k}) \mid 1 \leq i_{1} < \cdots < i_{k} \leq n\}$, then we have that $\dim \bigwedge^{k} T_{p}^{\ast}M = \binom{n}{k}$, and if $(U,\varphi) = (U,x^{1},\ldots,x^{n})$ is a chart on $M$, then the wedge products $\{(dx^{I})_{p} = (dx^{i_{1}})_{p} \wedge \cdots \wedge (dx^{i_{k}})_{p} \mid (i_{1},\ldots,i_{k}) \in \cI_{k,n}\}$ form a basis for $\bigwedge^{k}T_{p}^{\ast}M$. Thus, any $\omega_{p} \in A_{k}(T_{p}M)$ can be written as
\begin{align}
    \omega_{p} = \sum_{I \in \cI_{k,n}} a_{I}(p)(dx^{I})_{p}
\end{align}
where $a_{I}(p)$'s are real numbers. More generally, any $k$-form $\omega$ on $U$ can be written as
\begin{align}
    \omega = \sum_{I \in \cI_{k,n}} a_{I} dx^{I}
\end{align}
where $a_{I}:U \to \R$ are real-valued functions on $U$. If $f_{1},\ldots,f_{k} \in C^{\infty}(U)$ are smooth functions on $U$, then we have
\begin{align}
    (df_{1} \wedge \cdots \wedge df_{k})\left( \frac{\partial}{\partial x^{i_{1}}}, \ldots, \frac{\partial}{\partial x^{i_{k}}} \right) = \det\left( \frac{\partial f_{j}}{\partial x^{i_{l}}} \right)_{1 \leq j,l \leq k}.
\end{align}
We also have the notion of bundles of $k$-forms on $M$. The case $k=1$ gives the cotangent bundle $T^{\ast}M$. These bundles are defined as
\begin{align}
    \bigwedge^{k}T^{\ast}M = \bigsqcup_{p \in M} \bigwedge^{k}T_{p}^{\ast}M = \bigsqcup_{p \in M} A_{k}(T_{p}M)
\end{align}
with projection $\pi:\bigwedge^{k}T^{\ast}M \to M$ defined as $\pi(\omega_{p}) = p$ for all $\omega_{p} \in \bigwedge^{k}T_{p}^{\ast}M$. A $k$-form on $M$, then, is a section of the bundle $\bigwedge^{k}T^{\ast}M$. To induce a topology on $\bigwedge^{k}T^{\ast}M$, we can use the charts on $M$ as we did with the tangent and cotangent bundles. Let $(U,\varphi) = (U,x^{1},\ldots,x^{n})$ be a chart on $M$. Also adopt the notation $\bigwedge^{k}T^{\ast}U = \pi^{-1}(U) = \bigsqcup_{p \in U} \bigwedge^{k}T_{p}^{\ast}M$. Then we have the bijective map
\begin{align}
    \tilde{\varphi}:\bigwedge^{k}T^{\ast}U \to \varphi(U) \times \R^{\binom{n}{k}}, \qquad \omega_{p} \mapsto (\varphi(p),a_{I_{1}}(p),\ldots,a_{I_{\binom{n}{k}}}(p)).
\end{align}
With the standard topology on $\varphi(U) \times \R^{\binom{n}{k}}$, we can give $\bigwedge^{k}T^{\ast}U$ the topology induced by $\tilde{\varphi}$. Then, letting $\cB_{U}$ be the collection of all open sets in $\bigwedge^{k}T^{\ast}U$, we can give $\bigwedge^{k}T^{\ast}M$ the topology generated by $\cB = \bigcup_{(U,\varphi)} \cB_{U}$, where the union is taken over all charts $(U,\varphi)$ on $M$. The exact same analysis as tangent bundles shows that $\bigwedge^{k}T^{\ast}M$ is a smooth manifold of dimension $n + \binom{n}{k}$, and the projection $\pi:\bigwedge^{k}T^{\ast}M \to M$ is a smooth map. The case for $k = 0$ is trivial, as $\bigwedge^{0}T^{\ast}M = M \times \R$ is a smooth manifold of dimension $n+1$, and the projection $\pi:\bigwedge^{0}T^{\ast}M \to M$ is the projection onto the first factor. A smooth $k$-form $\omega$ on $M$ is a smooth section of $\pi$. That is, $\omega$ is smooth if $\omega:M \to \bigwedge^{k}T^{\ast}M$ is a smooth map between smooth manifolds. Also, a $k$-form $\omega$ on $M$ is smooth if and only if for any chart $(U,\varphi) = (U,x^{1},\ldots,x^{n})$ on $M$, the functions $a_{I}:U \to \R$ defined as $\omega = \sum_{I \in \cI_{k,n}} a_{I} dx^{I}$ are smooth. Again, for $k = 0$, a smooth section of $\pi:M \times \R \to M$ is a smooth map $f:M \to M \times \R$ such that $\pi(f(p)) = p$ for all $p \in M$. That is, $f(p) = (p,g(p))$ for some smooth function $g:M \to \R$. Thus, a smooth 0-form on $M$ is simply a smooth function on $M$.
\\ \\
\textit{October 9th.}

Same as before, given a smooth map $F:N \to M$ between smooth manifolds $N$ and $M$, we can define the pullback of a $k$-form $\omega$ on $M$ by $F$ as the $k$-form $F^{\ast}\omega$ on $N$ defined as $(F^{\ast}\omega)_{p} = F^{\ast}_{p}(\omega_{F(p)})$ for all $p \in N$. That is, $(F^{\ast}\omega)_{p}:T_{p}N^{k} \to \R$ is the alternating $k$-linear map defined as
\begin{align}
    (F^{\ast}\omega)_{p}(X_{1},\ldots,X_{k}) = \omega_{F(p)}(F_{\ast,p}(X_{1}),\ldots,F_{\ast,p}(X_{k}))
\end{align}
for all $X_{1},\ldots,X_{k} \in T_{p}N$. The pullback of a $k$-form by a smooth map is also smooth. That is, if $\omega$ is a smooth $k$-form on $M$, then $F^{\ast}\omega$ is a smooth $k$-form on $N$. The proof is the same as the case for 1-forms.
\\ \\
From here, we define
\begin{align}
    \Omega^{k}(M) = \{\omega:M \to \bigwedge^{k}T^{\ast}M \mid \omega \text{ is smooth}\},
\end{align}
the vector space of all smooth $k$-forms on $M$, which is also a $C^{\infty}(M)$-module via $(f\omega)(p) = f(p)\omega_{p}$ for all $f \in C^{\infty}(M)$, $\omega \in \Omega^{k}(M)$, and $p \in M$.
\\ \\
\textit{October 12th.}

The wedge product of forms is also defined here, as follows. Let $\omega \in \Omega^{k}(M)$ and $\eta \in \Omega^{l}(M)$ be $k$- and $l$-forms on a smooth manifold $M$. Then we can define the wedge product $\omega \wedge \eta$ as the $(k+l)$-form on $M$ defined as
\begin{align}
    (\omega \wedge \eta)_{p} = \omega_{p} \wedge \eta_{p} = \frac{1}{k!l!} \sum_{\sigma \in S_{k+l}} \sgn(\sigma) \omega_{p}(X_{\sigma(1)},\ldots,X_{\sigma(k)}) \eta_{p}(X_{\sigma(k+1)},\ldots,X_{\sigma(k+l)})
\end{align}
for all $p \in M$ and $X_{1},\ldots,X_{k+l} \in T_{p}M$.
\begin{proposition}
    If $\omega \in \Omega^{k}(M)$ and $\eta \in \Omega^{l}(M)$ are smooth $k$- and $l$-forms on a smooth manifold $M$, then $\omega \wedge \eta$ is a smooth $(k+l)$-form on $M$.
\end{proposition}
\begin{proof}
    Let $(U,\varphi) = (U,x^{1},\ldots,x^{n})$ be a chart on $M$. Then we have
    \begin{align}
        \omega = \sum_{I \in \cI_{k,n}} a_{I} dx^{I}, \qquad \eta = \sum_{J \in \cI_{l,n}} b_{J} dx^{J}
    \end{align}
    for some smooth functions $a_{I}:U \to \R$ and $b_{J}:U \to \R$. Then we have
    \begin{align}
        (\omega \wedge \eta) = \sum_{I,J} a_{I}b_{J} dx^{I} \wedge dx^{J}.
    \end{align}
    Since the product of smooth functions is smooth, we have that each of the functions $a_{I}b_{J}$ are smooth. Thus, $\omega \wedge \eta$ is a smooth $(k+l)$-form on $M$.
\end{proof}
Again, the wedge product is $\R$-bilinear, $C^{\infty}(M)$-bilinear, associative, and anti-commutative. From here, for a smooth manifold $M$ of dimension $n$, we define
\begin{align}
    \Omega^{\ast}(M) = \bigoplus_{k=0}^{n} \Omega^{k}(M) = \bigoplus_{k=0}^{\infty} \Omega^{k}(M),
\end{align}
which is a graded $\R$-algebra with respect to the wedge product. A smooth map $F:N \to M$ between smooth manifolds $N$ and $M$, then, induces a graded $\R$-algebra homomorphism $F^{\ast}:\Omega^{\ast}(M) \to \Omega^{\ast}(N)$. We first show the pullback for general $k$-forms. For a $k$-form $\omega$ on $M$, we have
\begin{align}
    (F^{\ast}\omega)_{p}(X_{1},\ldots,X_{k}) = \omega_{F(p)}(F_{\ast,p}(X_{1}),\ldots,F_{\ast,p}(X_{k}))
\end{align}
for all $p \in N$ and $X_{1},\ldots,X_{k} \in T_{p}N$. Then we have
\begin{align}
    (F^{\ast}(\omega \wedge \eta))_{p}(X_{1},\ldots,X_{k+l}) &= (\omega \wedge \eta)_{F(p)}(F_{\ast,p}(X_{1}),\ldots,F_{\ast,p}(X_{k+l})) \notag \\ &= \frac{1}{k!l!} \sum_{\sigma \in S_{k+l}} \sgn(\sigma) \omega_{F(p)}(F_{\ast,p}(X_{\sigma(1)}),\ldots,F_{\ast,p}(X_{\sigma(k)})) \notag \\ &\qquad \cdot \eta_{F(p)}(F_{\ast,p}(X_{\sigma(k+1)}),\ldots,F_{\ast,p}(X_{\sigma(k+l)})) \notag \\ &= (F^{\ast}\omega \wedge F^{\ast}\eta)_{p}(X_{1},\ldots,X_{k+l})
\end{align}
and
\begin{align}
    (F^{\ast}(f\omega+g\eta))_{p}(X_{1},\ldots,X_{k}) &= (f\omega+g\eta)_{F(p)}(F_{\ast,p}(X_{1}),\ldots,F_{\ast,p}(X_{k})) \notag \\ &= f(F(p))\omega_{F(p)}(F_{\ast,p}(X_{1}),\ldots,F_{\ast,p}(X_{k})) + g(F(p))\eta_{F(p)}(F_{\ast,p}(X_{1}),\ldots,F_{\ast,p}(X_{k})) \notag \\ &= (F^{\ast}f)(p)(F^{\ast}\omega)_{p}(X_{1},\ldots,X_{k}) + (F^{\ast}g)(p)(F^{\ast}\eta)_{p}(X_{1},\ldots,X_{k}) \notag \\ &= ((F^{\ast}f)(F^{\ast}\omega) + (F^{\ast}g)(F^{\ast}\eta))_{p}(X_{1},\ldots,X_{k})
\end{align}
showing that $F^{\ast}(\omega \wedge \eta) = F^{\ast}\omega \wedge F^{\ast}\eta$ and $F^{\ast}(f\omega+g\eta) = (F^{\ast}f)(F^{\ast}\omega) + (F^{\ast}g)(F^{\ast}\eta)$ for all $\omega,\eta \in \Omega^{\ast}(M)$ and $f,g \in C^{\infty}(M)$. Now, we show that if $\omega \in \Omega^{k}(M)$, then $F^{\ast}\omega \in \Omega^{k}(N)$. Let $(U,\varphi) = (U,x^{1},\ldots,x^{n})$ be a chart on $M$, and let $(V,\psi) = (V,y^{1},\ldots,y^{m})$ be a chart on $N$, with $F(U) \subseteq V$. Then $F^{\ast}\omega$ is smooth since
\begin{align}
    (F^{\ast}\omega)_{p} = F^{\ast}_{p}(\omega_{F(p)}) &= \sum_{I \in \cI_{k,n}} a_{I}(F(p)) F^{\ast}_{p}(dx^{I}) = \sum_{I \in \cI_{k,n}} a_{I}(F(p)) d(x^{i_{1}} \circ F)_{p} \wedge \cdots \wedge d(x^{i_{k}} \circ F)_{p} \notag \\
    &= \sum_{I \in \cI_{k,n}} a_{I}(F(p)) d(F^{\ast}x^{i_{1}})_{p} \wedge \cdots \wedge d(F^{\ast}x^{i_{k}})_{p}
\end{align}
where each function $a_{I} \circ F$ and each function $x^{i} \circ F$ is smooth on $F^{-1}(U)$. Hence, $F^{\ast}\omega$ is a smooth $k$-form on $N$.

\subsection{Exterior Differentiation}

Let $A = \bigoplus_{n=0}^{\infty} A_{n}$ be a graded algebra over $\R$. An $\R$-linear map $D:A \to A$ is said to be an antiderivation of degree $+1$ if $D(ab) = (Da)b + (-1)^{\deg a}a(Db)$ for all $a \in A_{k}$ and $b \in A_{l}$, and $Da \in A_{k+1}$ whenever $a \in A_{k}$. In particular, let $A = \Omega^{\ast}(M)$ be the graded algebra of smooth differential forms on a smooth manifold $M$, and define the exterior derivative $d:\Omega^{\ast}(M) \to \Omega^{\ast}(M)$ as the $\R$-linear map defined as
\begin{align}
    d\left(\sum a_{I} dx^{I}\right) = \sum da_{I} \wedge dx^{I}
\end{align}
after coordinates are chosen. Then $d$ is an antiderivation of degree $+1$, satisfies $d^{2} = 0$, and for any $f \in C^{\infty}(M)$, $df \in \Omega^{1}(M)$ is the differential of $f$. Moreover, if $D$ is any antiderivation of degree $+1$ on $\Omega^{\ast}(M)$ such that $D^{2} = 0$ and $Df = df$ for all $f \in C^{\infty}(M)$, then $D = d$. That is, the exterior derivative is the unique antiderivation of degree $+1$ on $\Omega^{\ast}(M)$ that squares to zero and agrees with the differential on smooth functions.

\section{Orientable Manifolds}
\textit{October 14th.}

We first provide the theory for a finite dimensional vector space $V$ over $\R$, of dimension $n$. On the collection of all ordered bases of $V$, we define an equivalence relation as $(v_{1},\ldots,v_{n}) \sim (w_{1},\ldots,w_{n})$ if and only if the change of basis matrix from $(v_{1},\ldots,v_{n})$ to $(w_{1},\ldots,w_{n})$ has a positive determinant. Reflexivity, symmetry, and transitivity may be easily verified. This equivalence relation on the set of ordered basis of $V$ partitions it into 2 equivalence classes. An orientation on $V$ is a choice between these 2 equivalence clases or, equivalently, a choice of an ordered basis and its corresponding equivalence class. We say $V$ is oriented after making such a choice. For a finite dimensional vector space $V$ over $\R$, we shall use $\mu$ and $-\mu$ to denote the two orientations.

Lifting this theory, let $M$ be a smooth manifold of dimension $n$. Recall that a vector field $X$ is a section of the tangent bundle $TM$ such that $X:M \to TM$ satisfies $\pi \circ X = \id_{M}$, and a(n ordered) frame $(X_{1},X_{2},\ldots,X_{n})$ on $M$ is a collection of vector fields on $M$ such that for all $p \in M$, $(X_{1,p},X_{2,p},\ldots,X_{n,p})$ is a(n ordered) basis for $T_{p}M$.

\begin{definition}
    A \eax{pointwise orientation} $\mu$ on a smooth manifold $M$ is a choice of orientation $\mu_{p}$ on $T_{p}M$ for all $p \in M$. It is said to be continuous at a point $p \in M$ if $p$ has a neighbourhood $U$ on which $\mu$ is represented by a continuous frame on $U$; that is, there exist continuous vector fields $X_{1},\ldots,X_{n}$ on $U$ such that for all $q \in U$, $[(X_{1,p},X_{2,p},\ldots,X_{n,p})] = \mu_{p}$.

    An \eax{orientation} on $M$ is a pointwise orientation $\mu$ that is continuous for all $p \in M$. A manifold on which there exists an orientation is called an \eax{orientable manifold}. A manifold with a choice of orientation is called a \eax{oriented manifold}.
\end{definition}

\begin{proposition}
    Let $M$ be a smooth orientable connected manifold. Then $M$ has exactly 2 orientations.
\end{proposition}
\begin{proof}
    Take two orientations $\mu,\nu$ on $M$. Define the map $f:M \to \{\pm 1\}$ as $f(p) = 1$ if $\mu_{p} = \nu_{p}$ and $f(p) = -1$ if $\mu_{p} \neq \nu_{p}$. We claim that $f$ is a continuous map. Let $p \in U$, where $U$ is a connected neighbourhood of $p$. Then continuity at $p$ gives $\mu = (X_{1},\ldots,X_{n})$ and $\nu = (Y_{1},\ldots,Y_{n})$, two collections of continuous vector fields on $U$ which form a frame. Thus, there exists a matrix of continuous functions $A:U \to M_{n}(\R)$ such that
    \begin{align}
        (Y_{1,q},\ldots,Y_{n,q}) = (X_{1,q},\ldots,X_{n,q})A(q) \qquad \text{ for all } q \in U.
    \end{align}
    This implies that the map $\det(A(\cdot)):U \to \R^{\times}$ is a continuous function; so it must be positive on the entirety of $U$, or negative. Moreover, $f:U \to \{\pm 1\}$ may be written as $f(q) = \sgn(\det(A(q)))$ for all $q \in U$, showing that $f$ is locally constant. Since $M$ is connected, $f$ is constant, finally implying that $\mu = \nu$ or $\mu = -\nu$.
\end{proof}

\textit{October 16th.}

\begin{proposition}
    Let $M$ be a smooth manifold of dimension $n$. A frame $(X_{1},\ldots,X_{n})$ on $M$ gives a pointwise continuous orientation on $M$ if and only if for all $p \in M$, there exists a chart $(U,\varphi) = (U,x^{1},\ldots,x^{n})$ on $M$ with $p \in U$ such that $(dx^{1} \wedge \cdots \wedge dx^{n})_{p}(X_{1,p},\ldots,X_{n,p}) > 0$.
\end{proposition}
\begin{proof}
    Let $(X_{1},\ldots,X_{n})$ be a frame on $M$ which gives a pointwise continuous orientation on $M$. Let $p \in M$. Then there exists a neighbourhood $U$ of $p$ and a continuous frame $(Y_{1},\ldots,Y_{n})$ on $U$ such that for all $q \in U$, $[(Y_{1,q},\ldots,Y_{n,q})] = [(X_{1,q},\ldots,X_{n,q})]$. $[(\frac{d}{dx^{1}})_{q},\ldots,(\frac{d}{dx^{n}})_{q}]$ is an orientation on $T_{q}M$ for all $q \in U$. Since $Y_{1},\ldots,Y_{n}$ are continuous vector fields on $U$, the coordinate functions $a_{ij}:U \to \R$ defined as $Y_{i} = \sum_{j=1}^{n} a_{ij}(\frac{d}{dx^{j}})$ are continuous. Thus, the map $\det(a_{ij}):U \to \R^{\times}$ is continuous, and so it must be positive on the entirety of $U$, or negative. If this map is negative, we can replace $Y_{1}$ with $-Y_{1}$ to make it positive. Then we have
    \begin{align}
        (dx^{1} \wedge \cdots \wedge dx^{n})_{q}(Y_{1,q},\ldots,Y_{n,q}) = \det(a_{ij}(q)) > 0.
    \end{align}
    Finally, since $[(Y_{1,q},\ldots,Y_{n,q})] = [(X_{1,q},\ldots,X_{n,q})]$, we have that
    \begin{align}
        (dx^{1} \wedge \cdots \wedge dx^{n})_{q}(X_{1,q},\ldots,X_{n,q}) > 0.
    \end{align}
    For the converse, let $(X_{1},\ldots,X_{n})$ be a frame on $M$ such that for all $p \in M$, there exists a chart $(U,\varphi) = (U,x^{1},\ldots,x^{n})$ on $M$ with $p \in U$ such that $(dx^{1} \wedge \cdots \wedge dx^{n})_{p}(X_{1,p},\ldots,X_{n,p}) > 0$. Then for all $q \in U$, we have
    \begin{align}
        (dx^{1} \wedge \cdots \wedge dx^{n})_{q}(X_{1,q},\ldots,X_{n,q}) > 0.
    \end{align}
    We show that $[(X_{1,q},\ldots,X_{n,q})] = [(\frac{d}{dx^{1}})_{q},\ldots,(\frac{d}{dx^{n}})_{q}]$ for all $q \in U$. Let $A(q)$ be the change of basis matrix from $(X_{1,q},\ldots,X_{n,q})$ to $((\frac{d}{dx^{1}})_{q},\ldots,(\frac{d}{dx^{n}})_{q})$. Then we have
    \begin{align}
        (dx^{1} \wedge \cdots \wedge dx^{n})_{q}(X_{1,q},\ldots,X_{n,q}) = \det(A(q)) > 0.
    \end{align}
    This shows that the frame $(X_{1,q},\ldots,X_{n,q})$ is positively oriented with respect to the coordinate frame $((\frac{d}{dx^{1}})_{q},\ldots,(\frac{d}{dx^{n}})_{q})$ for all $q \in U$, and hence the frame $(X_{1},\ldots,X_{n})$ gives a pointwise continuous orientation on $M$.
\end{proof}

\begin{proposition}
    Let $M$ be a smooth manifold of dimension $n$. Then $M$ is orientable if and only if there exists a smooth nowhere vanishing $n$-form $\omega \in \Omega^{n}(M)$.
\end{proposition}
\begin{proof}
    Suppose $M$ is orientable, and $X_{1},\ldots,X_{n}$ is a frame on $M$ which gives a pointwise continuous orientation on $M$. Then for all $p \in M$, there exists a chart $(U,\varphi) = (U,x^{1},\ldots,x^{n})$ on $M$ with $p \in U$ such that $(dx^{1} \wedge \cdots \wedge dx^{n})_{p}(X_{1,p},\ldots,X_{n,p}) > 0$, or on $U$, $(dx^{1} \wedge \cdots \wedge dx^{n})(X_{1},\ldots,X_{n}) > 0$. Vary $p \in M$ to get an open cover $\{U_{\alpha}\}$ of $M$ with charts $(U_{\alpha},\varphi_{\alpha}) = (U_{\alpha},x^{1}_{\alpha},\ldots,x^{n}_{\alpha})$ on $M$ such that $(dx^{1}_{\alpha} \wedge \cdots \wedge dx^{n}_{\alpha})(X_{1},\ldots,X_{n}) > 0$ on $U_{\alpha}$. Let $\{\rho_{\alpha}\}$ be a smooth partition of unity subordinate to the open cover $\{U_{\alpha}\}$. Then we can define a smooth $n$-form $\omega \in \Omega^{n}(M)$ as
    \begin{align}
        \omega = \sum_{\alpha} \rho_{\alpha} dx^{1}_{\alpha} \wedge \cdots \wedge dx^{n}_{\alpha}.
    \end{align}
    Then for all $p \in M$, we have
    \begin{align}
        \omega_p(X_{1,p},\ldots,X_{n,p}) = \sum_{\alpha} \rho_{\alpha}(p) (dx^{1}_{\alpha} \wedge \cdots \wedge dx^{n}_{\alpha})_p(X_{1,p},\ldots,X_{n,p}) > 0.
    \end{align}
    This shows that $\omega$ is a smooth nowhere vanishing $n$-form on $M$. For the converse, suppose there exists a smooth nowhere vanishing $n$-form $\omega \in \Omega^{n}(M)$. Then for all $p \in M$, we have $\omega_{p} \neq 0$, and hence there exists a basis $(X_{1,p},\ldots,X_{n,p})$ for $T_{p}M$ such that $\omega_{p}(X_{1,p},\ldots,X_{n,p}) > 0$. Vary $p \in M$ to get a frame $(X_{1},\ldots,X_{n})$ on $M$ such that $\omega(X_{1},\ldots,X_{n}) > 0$. 
    Define an orientation on $T_{p}M$ by declaring an ordered basis $(v_{1},\ldots,v_{n})$ to be positive if and only if
    \begin{align}
        \omega_{p}(v_{1},\ldots,v_{n}) > 0.
    \end{align}
    This is well-defined, since replacing a basis by another one multiplies the value of $\omega_{p}$ by the determinant of the change-of-basis matrix. It remains to verify continuity. Let $p \in M$ and choose a chart $(U,x^{1},\ldots,x^{n})$ around $p$. Since $\omega$ is nowhere vanishing, we can write
    \begin{align}
        \omega = f\,dx^{1} \wedge \cdots \wedge dx^{n}
    \end{align}
    on $U$, where $f:U \to \mathbb{R}$ is smooth and nowhere zero. After shrinking $U$ if necessary, $f$ has constant sign. Consequently, the coordinate frame $(\frac{\partial}{\partial x^{1}},\ldots,\frac{\partial}{\partial x^{n}})$, or the same frame with one vector multiplied by $-1$, represents the orientation defined above throughout $U$. Thus the pointwise orientation is continuous, so $M$ is orientable.
\end{proof}

\textit{October 21st.}

\begin{definition}
    Let $\{(U_{\alpha},\varphi_{\alpha}) = (U_{\alpha},x^{1}_{\alpha},\ldots,x^{n}_{\alpha})\}_{\alpha \in \Lambda}$ be an atlas on a smooth manifold $M$ of dimension $n$. It is termed an \eax{orientable atlas} if for all $\alpha,\beta \in \Lambda$ with $U_{\alpha} \cap U_{\beta} \neq \emptyset$, the Jacobian determinant of the transition map $\varphi_{\beta} \circ \varphi_{\alpha}^{-1}$ is positive on $\varphi_{\alpha}(U_{\alpha} \cap U_{\beta})$; that is,
    \begin{align}
        \det\left( \frac{\partial x^{i}_{\beta}}{\partial x^{j}_{\alpha}} \right)_{1 \leq i,j \leq n} > 0 \text{ on } \varphi_{\alpha}(U_{\alpha} \cap U_{\beta}).
    \end{align}
\end{definition}

\begin{proposition}
    Let $M$ be a smooth manifold of dimension $n$. Then $M$ is orientable if and only if there exists an orientable atlas on $M$.
\end{proposition}
\begin{proof}
    If $M$ is orientable, then there exists a frame $(X_{1},\ldots,X_{n})$ on $M$ such that there exist charts $(U,\varphi) = (U,x^{1},\ldots,x^{n})$ covering $M$ satisfying $(dx^{1} \wedge \cdots \wedge dx^{n})(X_{1},\ldots,X_{n}) > 0$. Let $\{(U_{\alpha},\varphi_{\alpha}) = (U_{\alpha},x^{1}_{\alpha},\ldots,x^{n}_{\alpha})\}_{\alpha \in \Lambda}$ be this atlas of charts on $M$. Then for all $\alpha,\beta \in \Lambda$ with $U_{\alpha} \cap U_{\beta} \neq \emptyset$, we have
    \begin{align}
         \det\left( \frac{\partial x^{i}_{\beta}}{\partial x^{j}_{\alpha}} \right) \cdot (dx^{1}_{\alpha} \wedge \cdots \wedge dx^{n}_{\alpha})(X_{1},\ldots,X_{n}) = (dx^{1}_{\beta} \wedge \cdots \wedge dx^{n}_{\beta})(X_{1},\ldots,X_{n}) > 0
    \end{align}
    implying that the Jacobian determinant of the transition map $\varphi_{\beta} \circ \varphi_{\alpha}^{-1}$ is positive on $\varphi_{\alpha}(U_{\alpha} \cap U_{\beta})$. Thus, $\{(U_{\alpha},\varphi_{\alpha})\}_{\alpha \in \Lambda}$ is an orientable atlas on $M$. Conversely, if there exists an orientable atlas $\{(U_{\alpha},\varphi_{\alpha})\}_{\alpha \in \Lambda}$ on $M$, let $\mu$ be the pointwise orientation on $M$ defined as follows. For all $p \in M$, choose $\alpha \in \Lambda$ such that $p \in U_{\alpha}$, and define $\mu_{p} = [((\frac{\partial}{\partial x^{1}_{\alpha}})_{p},\ldots,(\frac{\partial}{\partial x^{n}_{\alpha}})_{p})]$. This is well-defined, since if $p \in U_{\alpha} \cap U_{\beta}$, then the Jacobian determinant of the transition map $\varphi_{\beta} \circ \varphi_{\alpha}^{-1}$ is positive on $\varphi_{\alpha}(U_{\alpha} \cap U_{\beta})$, and hence
    \begin{align}
        \left[\left(\left(\frac{\partial}{\partial x^{1}_{\alpha}}\right)_{p},\ldots,\left(\frac{\partial}{\partial x^{n}_{\alpha}}\right)_{p}\right)\right] = \left[\left(\left(\frac{\partial}{\partial x^{1}_{\beta}}\right)_{p},\ldots,\left(\frac{\partial}{\partial x^{n}_{\beta}}\right)_{p}\right)\right].
    \end{align}
\end{proof}

Let $M$ be a connected smooth manifold of dimension $n$. Then $M$ is orientable if and only if there exists a smooth nowhere vanishing $n$-form $\omega \in \Omega^{n}(M)$; on the set of all smooth nowhere vanishing $n$-forms on $M$, we define an equivalence relation as $\omega \sim \eta$ if and only if $\omega = f\eta$ for some smooth function $f:M \to \R$ such that $f(p) > 0$ for all $p \in M$. This equivalence relation partitions the set of all smooth nowhere vanishing $n$-forms on $M$ into 2 equivalence classes. Thus, an orientation determined by a smooth nowhere vanishing $n$-form $\omega$ on $M$ is exactly the orientation given by a frame $(X_{1},\ldots,X_{n})$ on $M$ such that $\omega(X_{1},\ldots,X_{n}) > 0$. In particular, if $\omega$ and $\eta$ are smooth nowhere vanishing $n$-forms on $M$, then they determine the same orientation on $M$ if and only if $\omega = f\eta$ for some smooth function $f > 0$.

\subsection{Manifolds with Boundary}
\textit{October 24th.}

Hereforth, we use $\sH^{n} = \{(x^{1},\ldots,x^{n}) \in \R^{n} \mid x^{n} \geq 0\}$ to denote the closed upper half-space in $\R^{n}$.

\begin{definition}
    A topological space $X$ is \eax{locally $\sH^{n}$} if for all $p \in X$, there exists a neighbourhood $U$ of $p$ and a homeomorphism $\varphi:U \to V$ onto an open subset $V$ of $\sH^{n}$ such that $\varphi(p) = 0$. $M$ is termed a \eax{topological manifold with boundary} of dimension $n$ if it is Hausdorff, second countable, and locally $\sH^{n}$.
\end{definition}

A smooth atlas on $M$ is a collection of charts $\{(U_{\alpha},\varphi_{\alpha})\}_{\alpha \in \Lambda}$ such that for all $\alpha,\beta \in \Lambda$ with $U_{\alpha} \cap U_{\beta} \neq \emptyset$, the transition map $\varphi_{\beta} \circ \varphi_{\alpha}^{-1}$ is smooth on $\varphi_{\alpha}(U_{\alpha} \cap U_{\beta})$. A smooth manifold with boundary is a topological manifold with boundary equipped with a smooth atlas. 

For the closed half space $\sH^{n}$, we define
\begin{align}
    \partial \sH^{n} = \{(x^{1},\ldots,x^{n}) \in \R^{n} \mid x^{n} = 0\}, \qquad (\sH^{n})^{\circ} = \{(x^{1},\ldots,x^{n}) \in \R^{n} \mid x^{n} > 0\}
\end{align}
so that $\sH^{n} = (\sH^{n})^{\circ} \sqcup \partial \sH^{n}$. Note that if $U$ and $V$ are open subsets of $\sH^{n}$ and $f:U \to V$ is a diffeomorphism, then $f(\partial \sH^{n} \cap U) \subseteq \partial \sH^{n} \cap V$ and $f((\sH^{n})^{\circ} \cap U) \subseteq (\sH^{n})^{\circ} \cap V$. Here, $\partial \sH^{n}$ denotes the boundary of $\sH^{n}$, and $(\sH^{n})^{\circ}$ denotes the interior of $\sH^{n}$. One can show that $\partial \sH^{n}$ is a smooth manifold of dimension $n-1$, and $(\sH^{n})^{\circ}$ is a smooth manifold of dimension $n$. This can be extended to any smooth manifold with boundary $M$. We denote
\begin{align}
    \partial M = \{p \in M \mid \text{there exists a chart } (U,\varphi) \text{ on } M \text{ with } p \in U \text{ and } \varphi(p) \in \partial \sH^{n}\}, \qquad M^{\circ} = M \setminus \partial M.
\end{align}
One can show that $\partial M$ is a smooth manifold of dimension $n-1$, and $M^{\circ}$ is a smooth manifold of dimension $n$. We call $\partial M$ the \eax{boundary} of $M$, and $M^{\circ}$ the \eax{interior} of $M$. Note that if $M$ is a smooth manifold with boundary, then $\partial M$ is a smooth manifold without boundary.

Smooth functions on a smooth manifold with boundary $M$ are defined as follows. A function $f:M \to \R$ is smooth if for all $p \in M$, there exists a chart $(U,\varphi) = (U,x^{1},\ldots,x^{n})$ on $M$ with $p \in U$ such that $f \circ \varphi^{-1}$ is smooth on $\varphi(U)$. Smooth maps between smooth manifolds with boundary are defined similarly. Germs of smooth functions on a smooth manifold with boundary $M$ are defined just as for smooth manifolds without boundary. The tangent space $T_{p}M$ at a point $p \in M$ is defined as the set of derivations on the germs of smooth functions at $p$. We may then, analogously, define
\begin{align}
    T_{p}^{\ast}M = \Hom_{\R}(T_{p}M,\R), \qquad \bigwedge^{k}T_{p}^{\ast}M = A_{k}(T_{p}M)
\end{align}
and
\begin{align}
    TM = \bigsqcup_{p \in M} T_{p}M, \qquad T^{\ast}M = \bigsqcup_{p \in M} T_{p}^{\ast}M.
\end{align}
A smooth $k$-form on a smooth manifold with boundary $M$ is defined as a smooth section of the bundle $\bigwedge^{k}T^{\ast}M$. An orientation on a smooth manifold with boundary $M$, then, is one of three equivalent definitions: a smooth nowhere vanishing $n$-form $\omega \in \Omega^{n}(M)$, a pointwise continuous frame $(X_{1},\ldots,X_{n})$ on $M$, or an orientable atlas $\{(U_{\alpha},\varphi_{\alpha})\}_{\alpha \in \Lambda}$ on $M$.

Let $\beta \in \bigwedge^{k} V^{\ast} = A_{k}(V)$ be a $k$-covector on a finite dimensional vector space $V$ over $\R$, where $\dim V \geq k \geq 2$. If a vector $v \in V$ is fixed, then we can define a $(k-1)$-covector $i_{v}\beta \in \bigwedge^{k-1} V^{\ast}$ as
\begin{align}
    (i_{v}\beta)(v_{1},\ldots,v_{k-1}) = \beta(v,v_{1},\ldots,v_{k-1}).
\end{align}
If $k = 1$, then we define $i_{v}\beta \in \bigwedge^{0} V^{\ast} = \R$ as $(i_{v}\beta) = \beta(v)$. $k = 0$ admits $i_{v}\beta = 0$. Moreover, if $\beta^{1},\ldots,\beta^{k} \in V^{\ast}$ are $1$-covectors on $V$, then $\beta^{1} \wedge \cdots \wedge \beta^{k} \in \bigwedge^{k} V^{\ast}$ is a $k$-covector on $V$, and we have
\begin{align}
    i_{v}(\beta^{1} \wedge \cdots \wedge \beta^{k}) = \sum_{i=1}^{k} (-1)^{i-1} \beta^{1} \wedge \cdots \wedge \beta^{i-1} \wedge \beta^{i+1} \wedge \cdots \wedge \beta^{k}.
\end{align}
The same construction can be lifted to a smooth manifold, $M$. Let $\omega \in \Omega^{k}(M)$ be a smooth $k$-form on $M$, and let $X \in \mf{X}(M)$ be a vector field on $M$. Then we can define a smooth $(k-1)$-form $i_{X}\omega \in \Omega^{k-1}(M)$ as
\begin{align}
    (i_{X}\omega)_{p} = i_{X_{p}}(\omega_{p})
\end{align}
for all $p \in M$.

Note that if $M$ is an oriented smooth manifold with boundary, then $\partial M$ is an oriented smooth manifold without boundary.
\begin{definition}
    Let $M$ be a smooth manifold with boundary. Suppose $p \in \partial M$. A vector $X_{p} \in T_{p}M$ is said to be \eax{inward pointing} if
    \begin{enumerate}
        \item $X_{p} \notin T_{p}(\partial M)$, and
        \item there exists a smooth map $c:[0,\epsilon) \to M$ for some $\epsilon > 0$ such that $c(0) = p$, $c'(0) = X_{p}$, and $c((0,\epsilon)) \subseteq M^{\circ}$.
    \end{enumerate}
    It is said to be \eax{outward pointing} if $-X_{p}$ is inward pointing.
\end{definition}

\begin{definition}
    For a smooth manifold with boundary $M$, a vector field $X$ along the boundary $\partial M$ is a function $X:\partial M \to \bigsqcup_{p \in \partial M} T_{p}M$ such that $X(p) \in T_{p}M$ for all $p \in \partial M$. Such a vector field along the boundary is said to be smooth at $p \in \partial M$ if for any chart $(U,\varphi) = (U,x^{1},\ldots,x^{n})$ on $M$ with $p \in U$, $X$ when written as $X(q) = \sum_{i=1}^{n} a^{i}(q)(\frac{\partial}{\partial x^{i}})_{q}$ for $q \in U \cap \partial M$ has smooth coefficients $a^{i}:U \cap \partial M \to \R$. A vector field along the boundary is said to be smooth if it is smooth at all points $p \in \partial M$. It is outward pointing at $p \in \partial M$ if and only if $a^{n}(p) > 0$ in any chart $(U,\varphi) = (U,x^{1},\ldots,x^{n})$ on $M$ with $p \in U$. 
\end{definition}

\begin{proposition}
    For any smooth manifold with boundary, there exists a smooth outward pointing vector field along the boundary.
\end{proposition}
\begin{proof}
    Let $p \in \partial M$, and choose a chart $(U_{\alpha},x_{\alpha}^{1},\ldots,x_{\alpha}^{n})$ about $p$. Define $X_{\alpha} = -\frac{\partial}{\partial x_{\alpha}^{n}}$ on $U_{\alpha} \cap \partial M$. Then $X_{\alpha}$ is a smooth outward pointing vector field along the boundary on $U_{\alpha} \cap \partial M$. Let $\{(U_{\alpha},\varphi_{\alpha})\}_{\alpha \in \Lambda}$ be an open cover of $\partial M$ by such charts, and let $\{\rho_{\alpha}\}_{\alpha \in \Lambda}$ be a smooth partition of unity subordinate to this open cover. Then we can define a smooth outward pointing vector field along the boundary as
    \begin{align}
        X = \sum_{\alpha \in \Lambda} \rho_{\alpha} X_{\alpha}.
    \end{align}
    This is well-defined and smooth, and it is outward pointing at each point of $\partial M$.
\end{proof}

Now, let $M$ be an oriented smooth manifold with boundary of dimension $n$. Then $\partial M$ is an oriented smooth manifold without boundary of dimension $n-1$: let $\omega \in \Omega^{n}(M)$ be a smooth nowhere vanishing $n$-form on $M$ which determines the orientation on $M$. Let $X$ be a smooth outward pointing vector field along the boundary $\partial M$, via the above proposition. Then we can define a smooth nowhere vanishing $(n-1)$-form $\eta \in \Omega^{n-1}(\partial M)$ on $\partial M$ as
\begin{align}
    \eta = i_{X}\omega|_{\partial M}.
\end{align}
This $\eta$ cannot be zero; if it were at some point $p \in \partial M$, then $\omega_{p}(X_{p},v_{1},\ldots,v_{n-1}) = 0$ for all $v_{1},\ldots,v_{n-1} \in T_{p}(\partial M)$. But since $X_{p}$ is outward pointing, it is not in $T_{p}(\partial M)$, and hence $(X_{p},v_{1},\ldots,v_{n-1})$ forms a basis for $T_{p}M$. This contradicts the fact that $\omega$ is nowhere vanishing. Thus, $\eta$ is a smooth nowhere vanishing $(n-1)$-form on $\partial M$, and it determines an orientation on $\partial M$. Moreover, an ordered basis $(v_{1},v_{2},\ldots,v_{n-1})$ for $T_{p}(\partial M)$ is positively oriented with respect to the orientation on $\partial M$ if and only if $(X_{p},v_{1},v_{2},\ldots,v_{n-1})$ is positively oriented with respect to the orientation on $M$.